\documentclass[11pt]{article}

\usepackage[margin=1in]{geometry}
\usepackage{amsmath}
\usepackage{booktabs}
\usepackage{natbib}
\usepackage{hyperref}
\usepackage{graphicx}

\title{Intermittent Demand Forecasting for Automotive Retail:\\
An Operational Benchmark on Classical and Intermittent-Specific Methods}
\author{DealerIQ Project}
\date{\today}

\begin{document}
\maketitle

\begin{abstract}
We benchmark eight classical and intermittent-demand-specific forecasting methods
(Naive, Seasonal Naive, Simple Exponential Smoothing, automatic ETS, automatic ARIMA,
Croston's method, the Syntetos--Boylan Approximation, and Teunter--Syntetos--Babai)
against synthetically generated weekly vehicle-model demand spanning all four
Syntetos--Boylan--Croston demand-pattern classes (smooth, erratic, intermittent, lumpy),
using rolling-origin cross-validation across multiple independently-seeded datasets. We
report scale-free (MASE), percentage (sMAPE), and probabilistic (pinball loss, 80\%
interval coverage) accuracy per demand-pattern class, and test whether the
per-series-champion method's advantage over the mandatory Seasonal Naive floor is
statistically distinguishable from noise via a paired Wilcoxon signed-rank test. We then
describe how this benchmark's own forecasting pipeline is operationally integrated into
a real dealership-management system --- driving reorder-point computation and
stockout-risk alerting --- and discuss the honest limitations of a synthetic-data
benchmark for generalizing to real dealership demand.
\end{abstract}

\section{Introduction}

Automotive dealership parts and vehicle inventory exhibit demand patterns that range from
smooth, regularly-selling models to genuinely intermittent, slow-moving ones --- a
setting where classical forecasting methods tuned for smooth demand are known to
underperform \citep{croston1972forecasting, syntetos2005accuracy}. This work is the
benchmark and evaluation component of DealerIQ, a dealership intelligence platform whose
forecasting subsystem already selects, per (model, branch) series, a champion method from
a pool spanning both classical and intermittent-specific methods, and uses that
forecast's own predictive distribution to drive reorder-point computation and
stockout-risk alerting.

Two questions motivate this paper. First: across the full space of demand patterns a
real dealership's catalog spans, does the champion-selection process actually produce
better forecasts than the simplest reasonable baseline, and is that advantage real or an
artifact of a single lucky evaluation run? Second: how does forecast quality translate
into an operational decision (how much stock to hold, when to reorder) rather than
remaining an abstract accuracy number? We address the first with a multi-dataset,
statistically-tested benchmark; the second by describing the system's own, already-shipped
integration of forecast output into inventory decisions.

\section{Related Work}

The Syntetos--Boylan--Croston (SBC) framework classifies demand series by their average
inter-demand interval (ADI) and the squared coefficient of variation of nonzero demand
(CV\textsuperscript{2}) into smooth, erratic, intermittent, and lumpy classes
\citep{syntetos2005categorization}, motivating specialized methods for intermittent
demand originating with \citet{croston1972forecasting} and refined by the Syntetos--Boylan
Approximation \citep{syntetos2005accuracy} and Teunter--Syntetos--Babai
\citep{teunter2011intermittent}. \citet{hyndman2006another} introduced the Mean Absolute
Scaled Error (MASE) as a scale-free accuracy metric well-defined even for series with
long runs of zero demand, unlike percentage-error metrics; we use it as our primary
accuracy measure, following the same recommendation echoed in
\citet{hyndman2021forecasting}. \citet{makridakis2018statistical} survey the broader
tension between statistical and machine-learning forecasting methods and argue for
rigorous, reproducible comparison protocols, which this benchmark follows via
rolling-origin cross-validation. On the applied side, \citet{silver2016inventory} and
\citet{boylan2021intermittent} connect forecast quality to inventory-control decisions
(reorder points, safety stock) --- the same connection Section~\ref{sec:integration}
describes for this system specifically.

\section{Methodology}

\subsection{Synthetic data generation}

Weekly demand for each (vehicle model, branch) series is drawn from a Negative Binomial
generalized linear model with a log link: $\log \mu = \beta_0(\text{model}) + \text{trend}
+ \text{Fourier seasonal terms (annual)} + \text{promotion effects} + \text{branch
effect}$, with per-series dispersion tuned so the resulting series span all four SBC
classes. Each independently-seeded dataset generates 156 weeks (three years) of history
across the dealership's full simulated catalog. Full generator details are documented in
the project's own specification (\texttt{specs/004-synthetic-data/}) and are not repeated
here; this paper treats the generator as already-validated infrastructure and focuses on
the forecasting benchmark built on top of it.

\subsection{Method pool}

Eight methods, all via the \texttt{statsforecast} library: Naive; Seasonal Naive (the
mandatory champion-selection floor, described below); Simple Exponential Smoothing with
optimized smoothing parameter; automatic ETS (searches error/trend/seasonal component
combinations); automatic ARIMA (seasonal search space narrowed for tractable rolling-CV
runtime); Croston's classic method; the Croston--Syntetos--Boylan Approximation (SBA); and
Teunter--Syntetos--Babai (TSB, with commonly-cited default smoothing constants
$\alpha_d = \alpha_p = 0.1$, following worked examples in
\citet{teunter2011intermittent}).

\subsection{Rolling-origin cross-validation and champion selection}

Every method is evaluated on every series via rolling-origin (expanding-window)
cross-validation: a minimum 52-week training window, a 4-week step between folds, and a
12-week forecast horizon per fold. Per-series, per-method accuracy (MASE, sMAPE, pinball
loss at the 10th/50th/90th percentiles, and empirical 80\% interval coverage) is averaged
across folds. The champion method for a series is the one minimizing mean MASE
\emph{among methods that beat Seasonal Naive on MASE} --- a method is never selected as
champion unless it clears this floor, so Seasonal Naive itself becomes champion whenever
no other method genuinely improves on it. This selection rule, and the CV/metrics
implementation itself, are unchanged from the system's own operational forecasting
pipeline; this benchmark composes them without modification.

\subsection{Statistical comparison}

For each demand-pattern class (and pooled across all four), we extract the paired
$(\text{champion MASE}, \text{Seasonal Naive MASE})$ for every series across every
independently-seeded dataset and run a paired Wilcoxon signed-rank test, reporting the
test statistic, $p$-value, and the real sample size (number of paired series) alongside
every result. A class with too few series for a meaningful test is reported as such
rather than presented with false statistical confidence.

\section{Results}
\label{sec:results}

All numbers in this section are taken directly from a $K=3$ benchmark run
(three independently-seeded synthetic datasets, seeds 42--44) and are reproduced
verbatim in \texttt{docs/paper/results/leaderboard.md} and
\texttt{docs/paper/results/benchmark\_results.parquet}, which \texttt{make
benchmark-paper} regenerates from scratch.

\subsection{Dataset coverage}

Across the three seeded datasets, 198 (vehicle model, branch) series were generated
and classified, none discarded as degenerate. Table~\ref{tab:coverage} shows the
per-class series counts. Intermittent and lumpy demand dominate the generated
catalog (86 and 54 series respectively), with erratic the least represented (27
series) --- a property of the generator's own tuned dispersion parameters
(Section~\ref{sec:limitations}), not a claim about a real dealership's actual class
mix.

\begin{table}[htbp]
\centering
\begin{tabular}{lrrrrr}
\toprule
& Smooth & Erratic & Intermittent & Lumpy & Total \\
\midrule
$n$ series & 31 & 27 & 86 & 54 & 198 \\
\bottomrule
\end{tabular}
\caption{Series count by SBC class, pooled across all three seeded datasets ($K=3$).}
\label{tab:coverage}
\end{table}

\subsection{Accuracy by demand-pattern class}

Tables~\ref{tab:smooth}--\ref{tab:lumpy} report mean MASE, sMAPE, pinball loss (at
the 10th/50th/90th percentiles), and empirical 80\% interval coverage for every
method, per SBC class. Two consistent patterns hold across all four classes.
First, on MASE the mandatory Seasonal Naive floor and plain Naive are the worst or
second-worst performer in every class, confirming that the champion-selection
pool's non-trivial methods (AutoARIMA, AutoETS, SES, Croston/SBA, TSB) genuinely
improve on the floor rather than merely matching it. AutoARIMA is the best or
near-best method by MASE in the smooth, erratic, and lumpy classes (0.770, 0.763,
and 0.856 respectively); TSB and SES edge it out specifically on intermittent and
lumpy demand (TSB 0.911 on intermittent, SES 0.840 on lumpy), consistent with these
methods' intermittent-demand design intent.

Second, and disclosed honestly rather than smoothed over: the four methods with no
native parametric interval estimator (SES, Croston, SBA, TSB --- see
Section~\ref{sec:integration}'s discussion and this project's own ADR-025) achieve
only 59--63\% empirical 80\% coverage under the conformal-interval calibration
(\texttt{ConformalIntervals}, $n_{\text{windows}}=2$) used to give them intervals at
all, versus 78--97\% for the methods with native intervals (Naive, Seasonal Naive,
AutoETS, AutoARIMA). This is a real, measured undercoverage, not a hypothetical
concern: $n_{\text{windows}}=2$ is a deliberately disclosed default (ADR-025), and
this benchmark is the first measurement of its actual calibration quality. A larger
calibration window count is a legitimate future tuning knob this result motivates,
not a correctness bug to fix silently.

\begin{table}[htbp]
\centering\footnotesize
\begin{tabular}{lrrrrrr}
\toprule
Method & MASE & sMAPE & Pinball q10 & Pinball q50 & Pinball q90 & 80\% cov.\ \\
\midrule
AutoARIMA & 0.770 & 75.8\% & 0.22 & 0.64 & 0.35 & 81\% \\
AutoETS & 0.776 & 76.3\% & 0.23 & 0.65 & 0.35 & 82\% \\
TSB & 0.784 & 76.5\% & 0.29 & 0.65 & 0.42 & 60\% \\
CrostonClassic & 0.785 & 76.3\% & 0.30 & 0.65 & 0.42 & 60\% \\
CrostonSBA & 0.778 & 76.8\% & 0.29 & 0.65 & 0.43 & 60\% \\
SESOpt & 0.795 & 77.5\% & 0.30 & 0.66 & 0.42 & 60\% \\
SeasonalNaive & 1.004 & 97.7\% & 0.41 & 0.84 & 0.42 & 78\% \\
Naive & 1.035 & 99.2\% & 0.67 & 0.86 & 0.67 & 97\% \\
\bottomrule
\end{tabular}
\caption{Smooth class ($n = 31$ series).}
\label{tab:smooth}
\end{table}

\begin{table}[htbp]
\centering\footnotesize
\begin{tabular}{lrrrrrr}
\toprule
Method & MASE & sMAPE & Pinball q10 & Pinball q50 & Pinball q90 & 80\% cov.\ \\
\midrule
AutoARIMA & 0.763 & 91.1\% & 0.54 & 1.51 & 1.03 & 89\% \\
CrostonSBA & 0.766 & 91.4\% & 0.66 & 1.52 & 1.17 & 60\% \\
TSB & 0.779 & 91.7\% & 0.67 & 1.54 & 1.15 & 60\% \\
AutoETS & 0.774 & 91.5\% & 0.57 & 1.53 & 1.03 & 88\% \\
CrostonClassic & 0.780 & 91.6\% & 0.68 & 1.54 & 1.15 & 60\% \\
SESOpt & 0.780 & 91.8\% & 0.68 & 1.54 & 1.15 & 61\% \\
SeasonalNaive & 0.999 & 112.2\% & 1.16 & 1.98 & 1.19 & 79\% \\
Naive & 1.013 & 112.3\% & 1.87 & 1.97 & 1.87 & 96\% \\
\bottomrule
\end{tabular}
\caption{Erratic class ($n = 27$ series).}
\label{tab:erratic}
\end{table}

\begin{table}[htbp]
\centering\footnotesize
\begin{tabular}{lrrrrrr}
\toprule
Method & MASE & sMAPE & Pinball q10 & Pinball q50 & Pinball q90 & 80\% cov.\ \\
\midrule
TSB & 0.911 & 147.2\% & 0.13 & 0.34 & 0.24 & 61\% \\
SESOpt & 0.913 & 147.9\% & 0.14 & 0.34 & 0.24 & 61\% \\
AutoARIMA & 0.916 & 146.6\% & 0.11 & 0.34 & 0.20 & 88\% \\
AutoETS & 0.919 & 147.1\% & 0.12 & 0.34 & 0.21 & 88\% \\
CrostonSBA & 0.924 & 146.8\% & 0.12 & 0.34 & 0.24 & 62\% \\
CrostonClassic & 0.940 & 146.1\% & 0.12 & 0.35 & 0.24 & 62\% \\
SeasonalNaive & 1.005 & 84.4\% & 0.23 & 0.39 & 0.24 & 80\% \\
Naive & 1.024 & 85.8\% & 0.38 & 0.40 & 0.39 & 97\% \\
\bottomrule
\end{tabular}
\caption{Intermittent class ($n = 86$ series).}
\label{tab:intermittent}
\end{table}

\begin{table}[htbp]
\centering\footnotesize
\begin{tabular}{lrrrrrr}
\toprule
Method & MASE & sMAPE & Pinball q10 & Pinball q50 & Pinball q90 & 80\% cov.\ \\
\midrule
SESOpt & 0.840 & 149.2\% & 0.33 & 0.76 & 0.63 & 63\% \\
TSB & 0.846 & 148.6\% & 0.32 & 0.76 & 0.63 & 60\% \\
AutoARIMA & 0.856 & 145.8\% & 0.32 & 0.76 & 0.61 & 92\% \\
CrostonSBA & 0.883 & 147.8\% & 0.29 & 0.77 & 0.63 & 62\% \\
AutoETS & 0.885 & 148.2\% & 0.34 & 0.77 & 0.62 & 92\% \\
CrostonClassic & 0.902 & 147.6\% & 0.29 & 0.78 & 0.63 & 62\% \\
Naive & 0.917 & 78.2\% & 1.13 & 0.88 & 1.16 & 97\% \\
SeasonalNaive & 0.995 & 80.3\% & 0.70 & 0.94 & 0.72 & 85\% \\
\bottomrule
\end{tabular}
\caption{Lumpy class ($n = 54$ series).}
\label{tab:lumpy}
\end{table}

\subsection{Champion vs.\ Seasonal Naive: statistical significance}

Table~\ref{tab:wilcoxon} reports the paired Wilcoxon signed-rank test comparing
each series' champion-method MASE against its Seasonal Naive MASE, per SBC class
and pooled overall ($n = 198$ paired series). The champion beats Seasonal Naive on
88--100\% of series in every class, and the improvement is statistically
significant at $p < 0.0001$ in every class and overall --- this is not an artifact
of a single evaluation run: the same result holds pooled across three
independently-seeded datasets. Intermittent demand is the one class where the
champion does not win universally (88.4\% win rate, versus 100\% in smooth,
erratic, and lumpy), consistent with intermittent series' inherently higher
forecast-difficulty and the narrower absolute MASE gap observed in
Table~\ref{tab:intermittent} between the champion-eligible methods and the floor.

\begin{table}[htbp]
\centering
\begin{tabular}{lrrrr}
\toprule
SBC class & $n$ pairs & Champion win rate & Statistic & $p$-value \\
\midrule
Smooth & 31 & 100.0\% & 0.00 & $< 0.0001$ \\
Erratic & 27 & 100.0\% & 0.00 & $< 0.0001$ \\
Intermittent & 86 & 88.4\% & 0.00 & $< 0.0001$ \\
Lumpy & 54 & 100.0\% & 0.00 & $< 0.0001$ \\
\midrule
Overall & 198 & 94.9\% & 0.00 & $< 0.0001$ \\
\bottomrule
\end{tabular}
\caption{Paired Wilcoxon signed-rank test, champion MASE vs.\ Seasonal Naive MASE.}
\label{tab:wilcoxon}
\end{table}

\section{Operational Integration}
\label{sec:integration}

Forecast output is not an isolated accuracy exercise in this system --- it directly
drives two real, already-shipped operational behaviors.

\textbf{Reorder points.} The persisted champion forecast for a series already carries
three known points of each future week's predictive distribution: the point forecast
(treated as the median) and the 10th/90th percentiles of its 80\% prediction interval.
The reorder-point module uses these directly: for the default 90\% service level, the
reorder point is simply the 90th-percentile forecast summed over the supplier lead time;
for any other configured service level, it linearly interpolates between the three known
anchor quantiles rather than assuming a parametric distribution the champion method may
not actually produce. This means every reorder point in the system is a direct function
of the same forecast this benchmark evaluates --- a more accurate forecast (lower MASE,
better-calibrated coverage) translates immediately into a more accurate reorder point,
not just a better-looking accuracy number.

\textbf{Stockout alerting.} A separate, more urgent alert fires when current on-hand
stock is already below the champion forecast's \emph{point estimate} of demand over the
lead-time window --- distinct from the reorder-point signal, which uses the
quantile-inflated (safety-stock-including) demand estimate. This means even the
forecast's central tendency, not just its upper tail, is operationally load-bearing:
systematic bias in the point forecast (not just poor interval calibration) would
directly produce false or missed stockout alerts.

\section{Limitations}
\label{sec:limitations}

This benchmark is built entirely on synthetically-generated demand. The generator's
demand-generating process (Negative Binomial GLM with trend, seasonal, and promotional
terms) is a deliberately simplified model of real dealership demand, and the SBC-class
coverage it produces is a property of the generator's own tuned dispersion parameters,
not empirically measured from a real dealership's sales history. A method's relative
ranking on this synthetic benchmark is evidence about how it handles the \emph{specific
statistical properties} the generator produces (seasonality shape, promotion-effect
structure, the achieved ADI/CV\textsuperscript{2} distribution across series) --- it is
not, by itself, evidence that the same ranking would hold on a real dealership's demand,
which may have autocorrelation structure, regime changes, or demand-pattern mixes the
generator does not reproduce. We report the real sample size behind every statistical
claim in Section~\ref{sec:results} precisely so a reader can judge how much weight a
given result can bear, but no result here should be read as a claim about real-world
generalization without the caveat that only real dealership data (or, per the project's
own stated future protocol, an incorporated public retail dataset) could establish.

The champion-selection rule itself is also a specific design choice --- minimizing mean
MASE across CV folds, subject to beating Seasonal Naive --- and a different selection
criterion (e.g. optimizing pinball loss directly, or a class-specific rule) could produce
different champions and different aggregate results.

\section{Conclusion}

We built a reproducible, multi-dataset benchmark of eight forecasting methods across all
four Syntetos--Boylan--Croston demand-pattern classes, reusing this project's own
already-validated, operationally-deployed forecasting pipeline unchanged, and tested
whether champion-method selection produces a statistically real improvement over the
Seasonal Naive floor rather than an artifact of a single evaluation run. We connected
this benchmark to the system's real operational use of forecast output --- reorder-point
computation and stockout alerting --- so the accuracy numbers in Section~\ref{sec:results}
are not disconnected from what the forecasts are actually used for. The synthetic-data
limitation in Section~\ref{sec:limitations} bounds what this benchmark can claim; a
natural next step is incorporating real dealership or public retail sales data to test
whether these results hold outside the generator's own assumptions.

\bibliographystyle{plain}
\bibliography{references}

\end{document}
